\documentclass[aspectratio=169]{beamer}
\usepackage[utf8]{inputenc}
\usepackage[spanish]{babel}
\usepackage{graphicx}
\usepackage{listings}
\usepackage{xcolor}
\usepackage{hyperref}
\usepackage{tikz}

% Tema limpio y profesional
\usetheme{Madrid}
\usecolortheme{whale}
\setbeamertemplate{navigation symbols}{}

% Colores personalizados
\definecolor{mango}{RGB}{255, 140, 0}
\definecolor{verde}{RGB}{40, 167, 69}
\definecolor{rojo}{RGB}{220, 53, 69}
\definecolor{codigobg}{RGB}{248, 248, 248}

% Estilo para código
\lstset{
  basicstyle=\ttfamily\small,
  backgroundcolor=\color{codigobg},
  frame=single,
  framerule=0.5pt,
  rulecolor=\color{gray!30},
  breaklines=true,
  showstringspaces=false,
  aboveskip=8pt,
  belowskip=8pt,
  xleftmargin=8pt,
  xrightmargin=8pt
}

\title[Clasificador de Mangos con CNN]{Clasificador de Mangos con CNN \\ API, Docker \& Frontend}
\subtitle{Fases 6 y 7 — Backend + Frontend + QA}
\author{Abel Watt}
\institute{Proyecto de Aprendizaje Automático}
\date{\today}

\begin{document}

% ==================== PORTADA ====================
\begin{frame}
  \titlepage
\end{frame}

% ==================== ÍNDICE ====================
\begin{frame}{Contenido}
  \tableofcontents
\end{frame}

% ==================== SECCIÓN 1 ====================
\section{Resumen del Proyecto}

\begin{frame}{El problema}
  \begin{columns}
    \column{0.5\textwidth}
    \textbf{Objetivo:} Clasificar mangos en 5 tipos y determinar si son exportables.

    \vspace{1em}
    \textbf{Clases:}
    \begin{itemize}
      \item Tipo 1 — 100\% exportable
      \item Tipo 2 — 80\% exportable
      \item Tipo 3 — 60\% exportable
      \item Tipo 4 — 40\% NO exportable
      \item Tipo 5 — 20\% NO exportable
    \end{itemize}

    \column{0.5\textwidth}
    \textbf{Mis roles:}
    \begin{itemize}
      \item \textbf{Fase 6:} Backend / DevOps (API + Docker)
      \item \textbf{Fase 7:} Frontend / QA (Interfaz + pruebas)
    \end{itemize}

    \vspace{1em}
    \textbf{Tecnologías:}
    \begin{itemize}
      \item Flask + TensorFlow (backend)
      \item Docker + Gunicorn (despliegue)
      \item HTML5 + Bootstrap + JS (frontend)
      \item ResNet50 Transfer Learning (modelo)
    \end{itemize}
  \end{columns}
\end{frame}

\begin{frame}{Arquitectura del sistema}
  \centering
  \begin{tikzpicture}[node distance=2cm, auto]
    \node[draw, rounded corners, fill=orange!20] (user) {Usuario (navegador)};
    \node[draw, rounded corners, fill=blue!20, right=of user] (flask) {Flask API\\0.0.0.0:5000};
    \node[draw, rounded corners, fill=green!20, right=of flask] (tf) {TensorFlow\\ResNet50};
    \node[draw, rounded corners, fill=gray!20, below=of flask] (docker) {Docker\\Container};

    \draw[->] (user) -- node[above] {HTTP} (flask);
    \draw[->] (flask) -- node[above] {model.predict()} (tf);
    \draw[->] (tf) -- node[below] {JSON} (flask);
    \draw[->] (flask) -- node[below] {JSON} (user);
    \draw[dashed] (docker) -- (flask);
  \end{tikzpicture}

  \vspace{1em}
  \begin{center}
    \small\textbf{3 endpoints:} \texttt{GET /} $\cdot$ \texttt{POST /predict} $\cdot$ \texttt{GET /health}
  \end{center}
\end{frame}

% ==================== SECCIÓN 2 ====================
\section{Backend — Fase 6}

\begin{frame}{API Flask — api/app.py}
  \textbf{Decisiones de diseño:}
  \begin{itemize}
    \item Modelo cargado \textbf{una vez} al iniciar (no por request)
    \item Preprocesamiento automático según modelo (ResNet vs CNN propia)
    \item Sin archivos temporales en disco (procesamiento en memoria)
    \item Siguiendo el ejemplo \texttt{API\_COVID.py} del profesor
  \end{itemize}

  \vspace{0.5em}
  \textbf{Estructura del código:}
  \begin{enumerate}
    \item Configuración (\texttt{MODEL\_NAME}, clases, mapa exportable)
    \item Carga del modelo (\texttt{tf.keras.models.load\_model})
    \item \texttt{preparar\_imagen()} — preprocesamiento 224×224
    \item \texttt{inferir()} — predicción + formateo
    \item Endpoints Flask
  \end{enumerate}
\end{frame}

\begin{frame}[fragile]{Endpoint POST /predict}
  \begin{lstlisting}[language=Python]
@app.route("/predict", methods=["POST"])
def predict():
    if "image" not in request.files:
        return jsonify({"error": "..."}), 400
    file = request.files["image"]
    # Validar extensión, preprocesar, predecir...
    img_array = preparar_imagen(file)   # 224x224, normalizar
    resultado = inferir(img_array)      # model.predict()
    return jsonify(resultado)           # {tipo, exportable, probabilidad}
  \end{lstlisting}

  \vspace{0.5em}
  \textbf{Respuesta:}
  \begin{lstlisting}[language=JavaScript]
{"tipo": "Tipo_4", "exportable": false, "probabilidad": 94.47}
  \end{lstlisting}
\end{frame}

\begin{frame}{Manejo de errores}
  \centering
  \begin{tabular}{|c|c|l|}
    \hline
    \textbf{Caso} & \textbf{HTTP} & \textbf{Mensaje} \\
    \hline
    Sin imagen & 400 & \texttt{"No se envió ninguna imagen"} \\
    Nombre vacío & 400 & \texttt{"Nombre de archivo vacío"} \\
    Formato inválido & 400 & \texttt{"Formato no soportado: .txt"} \\
    Imagen corrupta & 500 & \texttt{"Error en la inferencia: ..."} \\
    Error interno & 500 & \texttt{"Error en la inferencia: ..."} \\
    \hline
  \end{tabular}

  \vspace{1em}
  Todos los errores se capturan con \texttt{try/except} y devuelven JSON estructurado.
\end{frame}

% ==================== SECCIÓN 3 ====================
\section{Docker — Despliegue}

\begin{frame}{Dockerfile}
  \begin{lstlisting}[language=Docker]
FROM python:3.10.4-slim
WORKDIR /app
RUN apt-get install -y libgl1-mesa-glx libglib2.0-0

COPY api/requirements.txt .
RUN pip install --no-cache-dir -r requirements.txt
COPY api/app.py api/
COPY frontend/ frontend/
COPY models/ models/

ENV TF_ENABLE_ONEDNN_OPTS=0
EXPOSE 5000
CMD ["gunicorn", "--bind", "0.0.0.0:5000",
     "--workers", "1", "--timeout", "600",
     "--max-requests", "1",
     "api.app:app"]
  \end{lstlisting}
\end{frame}

\begin{frame}{¿Por qué Gunicorn?}
  \begin{columns}
    \column{0.5\textwidth}
    \textbf{Flask dev server:}
    \begin{itemize}
      \item No es para producción
      \item Un solo hilo
      \item Sin manejo de timeout
      \item Advertencia explícita de Flask
    \end{itemize}

    \column{0.5\textwidth}
    \textbf{Gunicorn:}
    \begin{itemize}
      \item Servidor WSGI probado
      \item Workers gestionados
      \item Timeout configurable
      \item Reciclaje de workers
    \end{itemize}
  \end{columns}

  \vspace{1em}
  \textbf{Conflictos resueltos Gunicorn + TensorFlow en CPU:}

  \begin{tabular}{|c|c|}
    \hline
    \textbf{Problema} & \textbf{Solución} \\
    \hline
    oneDNN + fork & \texttt{TF\_ENABLE\_ONEDNN\_OPTS=0} \\
    Worker timeout & \texttt{--timeout 600} \\
    Memory leak & \texttt{--max-requests 1} \\
    RAM × workers & \texttt{--workers 1} \\
    \hline
  \end{tabular}
\end{frame}

\begin{frame}{Comandos Docker}
  \centering
  \begin{tabular}{|l|l|}
    \hline
    \textbf{Acción} & \textbf{Comando} \\
    \hline
    Build & \texttt{docker build -t mango-api .} \\
    Iniciar & \texttt{docker run -d -p 5000:5000 --name mango-api mango-api} \\
    Logs & \texttt{docker logs -f mango-api} \\
    Detener & \texttt{docker stop mango-api} \\
    Reiniciar & \texttt{docker restart mango-api} \\
    Borrar & \texttt{docker rm -f mango-api} \\
    \hline
  \end{tabular}

  \vspace{1em}
  \textbf{Tiempo de carga:} $\sim$90s en CPU (modelo ResNet50 de 165 MB)
\end{frame}

% ==================== SECCIÓN 4 ====================
\section{Frontend — Fase 7}

\begin{frame}[fragile]{Frontend HTML5 + Bootstrap + JS Vanilla}
  \textbf{Tecnologías (según guía):}
  \begin{itemize}
    \item HTML5 semántico
    \item Bootstrap 4.6 (CDN)
    \item JavaScript vanilla (\texttt{fetch API})
  \end{itemize}

  \vspace{0.5em}
  \textbf{Componentes:}

  \begin{tabular}{|c|l|}
    \hline
    \textbf{Componente} & \textbf{Implementación} \\
    \hline
    Input archivo & \texttt{<input type="file" accept=".jpg,.jpeg,.png,.bmp">} \\
    Preview & \texttt{FileReader.readAsDataURL()} \\
    Botón & \texttt{onclick="clasificar()"} \\
    Loader & Bootstrap spinner (importante: CPU $\sim$60s) \\
    Resultado & Badge verde/rojo + barra de progreso \% \\
    Errores & Alertas Bootstrap warning/danger \\
    \hline
  \end{tabular}
\end{frame}

\begin{frame}[fragile]{Flujo del Frontend}
  \begin{lstlisting}[language=JavaScript]
function clasificar() {
  const formData = new FormData();
  formData.append("image", input.files[0]);

  fetch("/predict", { method: "POST", body: formData })
    .then(res => res.ok ? res.json() : throwError(res))
    .then(data => renderResult(data))  // Card con tipo, exportable, %
    .catch(err => showError(err));
}
  \end{lstlisting}

  \vspace{0.5em}
  \textbf{Ruta relativa} \texttt{"/predict"} — funciona en desarrollo y producción sin cambios.
\end{frame}

% ==================== SECCIÓN 5 ====================
\section{Pruebas — QA}

\begin{frame}{Pruebas realizadas}
  \textbf{Pruebas funcionales (7.7):}

  \begin{tabular}{|l|c|c|}
    \hline
    \textbf{Prueba} & \textbf{Esperado} & \textbf{Resultado} \\
    \hline
    Imagen Tipo 4 & exportable: false, 94.47\% & \textcolor{verde}{OK} \\
    Health check & \texttt{\{"status":"ok"\}} & \textcolor{verde}{OK} \\
    Frontend HTML & HTTP 200 & \textcolor{verde}{OK} \\
    \hline
  \end{tabular}

  \vspace{1em}
  \textbf{Pruebas de borde (7.8):}

  \begin{tabular}{|l|c|c|}
    \hline
    \textbf{Prueba} & \textbf{Esperado} & \textbf{Resultado} \\
    \hline
    Sin archivo & 400 & \textcolor{verde}{OK} \\
    Archivo corrupto & 500 & \textcolor{verde}{OK} \\
    Formato inválido (.txt) & 400 & \textcolor{verde}{OK} \\
    \hline
  \end{tabular}
\end{frame}

% ==================== SECCIÓN 6 ====================
\section{Red y Conectividad}

\begin{frame}{Problema WSL2 + Tailscale}
  \textbf{Problema:} WSL2 tiene IP privada interna no accesible desde la red.

  \textbf{Solución — Port Proxy de Windows:}
  \begin{lstlisting}[language=Bash]
netsh interface portproxy add v4tov4
    listenport=5000 listenaddress=0.0.0.0
    connectport=5000 connectaddress=172.29.130.59
  \end{lstlisting}

  \vspace{0.5em}
  \begin{center}
  \begin{tikzpicture}[node distance=2.5cm]
    \node[draw, rounded corners] (laptop) {Laptop\\100.x.x.x};
    \node[draw, rounded corners, right=of laptop] (ts) {Tailscale\\100.82.22.10};
    \node[draw, rounded corners, right=of ts] (win) {Windows\\portproxy};
    \node[draw, rounded corners, right=of win] (wsl) {WSL2\\172.29.130.59:5000};
    \draw[->] (laptop) -- (ts);
    \draw[->] (ts) -- (win);
    \draw[->] (win) -- (wsl);
  \end{tikzpicture}
  \end{center}
\end{frame}

% ==================== SECCIÓN 7 ====================
\section{Demostración}

\begin{frame}{Demo — Resultado real}
  \centering
  \textbf{Entrada:} Imagen de mango Tipo 4 (no exportable)

  \vspace{1em}
  \texttt{curl -X POST -F "image=@mango.jpg" /predict}

  \vspace{1em}
  \begin{beamercolorbox}[wd=0.7\textwidth,center,rounded=true]{block body}
    \Large\texttt{\{"tipo": "Tipo\_4",}\\
    \texttt{\ \ "exportable": false,}\\
    \texttt{\ \ "probabilidad": 94.47\}}
  \end{beamercolorbox}

  \vspace{1em}
  \begin{itemize}
    \item Clasificación correcta: Tipo 4
    \item Exportable: false (correcto: 40\% no exportable)
    \item Confianza: 94.47\% (alta certeza)
  \end{itemize}
\end{frame}

% ==================== CIERRE ====================
\section{Conclusión}

\begin{frame}{Entregables completados}
  \begin{columns}
    \column{0.5\textwidth}
    \textbf{Fase 6 — API + Despliegue:}
    \begin{itemize}
      \item[\checkmark] \texttt{api/app.py} (3 endpoints)
      \item[\checkmark] \texttt{Dockerfile} (gunicorn)
      \item[\checkmark] \texttt{api/requirements.txt}
      \item[\checkmark] \texttt{docs/despliegue.md}
    \end{itemize}

    \column{0.5\textwidth}
    \textbf{Fase 7 — Frontend + QA:}
    \begin{itemize}
      \item[\checkmark] \texttt{frontend/index.html}
      \item[\checkmark] Input + preview + fetch
      \item[\checkmark] Manejo de errores
      \item[\checkmark] Pruebas funcionales y de borde
    \end{itemize}
  \end{columns}

  \vspace{2em}
  \begin{center}
    \Large\textbf{Stack:} Flask + TensorFlow + ResNet50 + Docker + Gunicorn + Bootstrap

    \vspace{1em}
    \normalsize Repositorio: \url{https://github.com/elable947/mango-cnn-classifier}
  \end{center}
\end{frame}

\begin{frame}{¿Preguntas?}
  \centering
  \Huge{🥭}

  \vspace{1em}
  \large Clasificador de Mangos con CNN

  \vspace{0.5em}
  \normalsize Fases 6 y 7 — Backend + Frontend + QA
\end{frame}

\end{document}
